\section{Premise exchange over the weak base}
\label{sec:q-exchange-ipra}

The two numerical calculi differ only in their sixth head.  Their
interderivability can be established in iPRA without using the
PA-refutation compiler.  This observation also prevents a completed
forward calibration for the original head from leaving a separate gap
for QAX11H.

For the numerical implication $I(a,b)=4p(a,b)$, iPRA proves
\[
 a\le I(a,b),\qquad b\le I(a,b).
\]
The only exceptional zero pairing has $a=b=0$, where both inequalities
still hold.  Consequently reflexivity and the two inference operations
\[
 u\le a\Longrightarrow u\le I(a,b),\qquad
 u\le b\Longrightarrow u\le I(a,b)
\]
are valid primitive recursive bound certificates.

The pinned file \file{M6/Logic/RestrictedImplication.lean} supplies a
finite, acyclic family of proof constructors from the following interface:
an implication constructor, MP, the restricted row
\[
 (((c\to u)\to a)\to((a\to c)\to(b\to c)))
 \quad\text{with }u\le a,
\]
and exactly the three bound operations just displayed.  Its constructors
derive $K$, Peirce, composition, exchange, and $S$.  They use fixed finite
compositions of previously defined constructors; they perform no recursion
on an arbitrary source formula and use neither a quantifier head nor an
equality/vacuity axiom.  The adapter
\file{PA196/Logic/ImplicationBasis.lean} identifies this restricted row
with the numerical $\rho_8$ used here\cite{q-basis}.

Interpret each of these constructors as a function producing a numerical
derivation list.  A row appends one axiom with its displayed parameters;
MP concatenates the two premise lists, shifts premise indices, and appends
one MP line.  Index shifting and concatenation have quantifier-free
correctness proofs by induction on the finite list length.  Every row's
side condition follows from the three arithmetic bound operations.
Therefore finite induction on this fixed family of definitions yields
primitive recursive list constructors whose validity iPRA proves,
uniformly in their numerical arguments.  This uses the displayed source
recipes, rather than inferring weak-base correctness from a generic
claim of PA-provability.

In particular there is a primitive recursive numerical proof constructor
$\operatorname{Exch}(a,b,c)$ with iPRA-provable conclusion
\[
 (a\to(b\to c))\to(b\to(a\to c)).                 \tag{Q-Ex}
\]
For additional finite evidence, the supplied original Q packet's
parametric $K,S,\mathrm{MP}$ certificate was replayed: it has 61 lines,
17 $K$ instances, 14 $S$ instances, and 30 MP steps, with no remaining
assumptions.  These $K,S$ instances are discharged by the preceding
restricted-row constructors.  The certificate is a symbolic formula
identity and hence is valid also at numerical zero aliases; it does not
treat numerical implication as an injective free constructor when
checking a numerical proof.

Let $P=U(D,A)$, $Q=U(D,I(A,H))$, and $R=U(D,H)$.
In $\mathcal C_4^Q$, the sixth-head axiom is $Q\to(P\to R)$.
Apply (Q-Ex) at $(Q,P,R)$ and MP to obtain the original sixth head
$P\to(Q\to R)$.  Conversely, (Q-Ex) at $(P,Q,R)$ turns the original
sixth head into the Q head.  Every other axiom and both inference rules
are unchanged.

Translate an arbitrary valid proof list by retaining unchanged axiom
lines, replacing each sixth-head leaf with this fixed proof block, and
maintaining a table from old line numbers to new conclusion indices.
An old MP or generalization line becomes the same rule using its mapped
premise indices.  This is primitive recursion on the input list length.
The cumulative invariant says that the translated prefix is valid and
that each recorded conclusion equals the corresponding old conclusion.
It is a bounded Boolean check on the two finite lists and their index
map, so its induction is quantifier-free.

\begin{theorem}[Uniform Q translations]
There are primitive recursive proof-list transformations $F_Q,G_Q$ such
that iPRA proves
\[
 \operatorname{Prf}_{\mathcal C_4}(e,a)
 \Longrightarrow\operatorname{Prf}_{\mathcal C_4^Q}(F_Q(e),a),
\]
\[
 \operatorname{Prf}_{\mathcal C_4^Q}(e,a)
 \Longrightarrow\operatorname{Prf}_{\mathcal C_4}(G_Q(e),a).
\]
\end{theorem}

Both transformations preserve the conclusion literally.  In particular,
the two refutation sentences are equivalent in iPRA.  The mathematical
ingredients are inherited and attributed; this section records their
weak-base justification and separates them from the harder PA calibration.
Exact copies, blob identities, and the replayed finite exchange certificate
are in \file{references/q_exchange/}.  No new proof-assistant build is
claimed for this composition.
